\fontsize{14pt}{22.5pt}\selectfont

\section{KẾT QUẢ THỰC NGHIỆM VÀ ĐÁNH GIÁ}
\label{ch:results}

\indent Chương này trình bày và phân tích kết quả thu được từ quá trình huấn luyện ba
thuật toán PPO, SAC và MA-POCA trên ba môi trường đã thiết kế. Nội dung được tổ chức
lần lượt theo từng môi trường, trong đó mỗi môi trường được minh họa bằng đường cong
phần thưởng, bảng số liệu và nhận xét đi kèm. Sau phần kết quả riêng của từng môi
trường, chương tiến hành so sánh tổng hợp theo bốn tiêu chí đánh giá đã nêu ở
Mục~\ref{sub:evalcriteria}, rồi kết thúc bằng phần thảo luận về khả năng ứng dụng thực
tiễn và những hạn chế còn tồn tại. Các số liệu thực đo được trích xuất trực tiếp từ
tệp sự kiện TensorBoard, trong khi các số liệu ước lượng cho những tổ hợp chưa huấn
luyện đầy đủ đều được ghi chú rõ ràng.

% ═══════════════════════════════════════════════
\subsection{Kết quả môi trường Football Table}

\indent Môi trường Football Table được huấn luyện chính bằng PPO kết hợp cơ chế tự
chơi. Trong cơ chế này, phần thưởng tuyệt đối của tác nhân có xu hướng dao động quanh
giá trị không, bởi khi cả hai phía cùng mạnh lên thì lợi thế của bên này bị bù trừ bởi
sự tiến bộ của bên kia. Vì vậy, chỉ số phản ánh tiến bộ tốt nhất trong cơ chế tự chơi
không phải là phần thưởng tuyệt đối mà là chỉ số xếp hạng. Trong thực nghiệm, chỉ số
xếp hạng của tác nhân PPO tăng ổn định từ giá trị khởi tạo một nghìn hai trăm và hội
tụ quanh mức một nghìn hai trăm mười chín sau khoảng hai trăm năm mươi nghìn bước, cho
thấy tác nhân liên tục cải thiện năng lực so với các phiên bản trước của chính mình.
Hình~\ref{fig:football_elo} biểu diễn tiến trình chỉ số xếp hạng thực đo này. Có thể
thấy chỉ số tăng gần như đơn điệu trong giai đoạn đầu rồi đi ngang, phản ánh trạng
thái cân bằng của quá trình tự chơi: khi cả hai phía đã đạt cùng một trình độ cao, lợi
thế của một bên khó còn tăng thêm và chỉ số xếp hạng ổn định lại.

\begin{figure}[H]
    \centering
    \begin{tikzpicture}
    \begin{axis}[
        width=13cm, height=6.5cm,
        xlabel={Số bước huấn luyện (triệu)},
        ylabel={Chỉ số xếp hạng ELO},
        xmin=0, xmax=0.5, ymin=1185, ymax=1225,
        legend pos=south east,
        grid=both, grid style={gray!20},
        tick label style={font=\small},
        label style={font=\small},
        scaled y ticks=false, yticklabel style={/pgf/number format/fixed},
    ]
    \addplot[ppocolor, thick, mark=*, mark size=1.5pt] coordinates {
        (0.05,1200.0)(0.10,1195.0)(0.15,1207.3)(0.20,1215.4)(0.25,1219.2)(0.30,1219.2)(0.35,1219.2)(0.40,1219.2)(0.45,1219.2)(0.50,1219.2)};
    \legend{PPO tự chơi (thực đo)}
    \end{axis}
    \end{tikzpicture}
    \caption{Tiến trình chỉ số xếp hạng ELO của PPO tự chơi trên Football Table}
    \label{fig:football_elo}
\end{figure}

\indent Để so sánh ba thuật toán trên cùng một thước đo, đề tài đánh giá cả ba bằng
phần thưởng trung bình khi đấu với cùng một đối thủ điều khiển bằng luật cố định.
Hình~\ref{fig:football_reward} biểu diễn đường cong phần thưởng đánh giá này, còn
Bảng~\ref{tab:football_results} tổng hợp các chỉ số định lượng tương ứng.

\begin{figure}[H]
    \centering
    \begin{tikzpicture}
    \begin{axis}[
        width=13cm, height=7cm,
        xlabel={Số bước huấn luyện (triệu)},
        ylabel={Phần thưởng đánh giá},
        xmin=0, xmax=0.5, ymin=-1, ymax=1,
        legend pos=south east,
        grid=both, grid style={gray!20},
        tick label style={font=\small},
        label style={font=\small},
    ]
    \addplot[ppocolor, thick, mark=none] coordinates {
        (0,-0.90)(0.05,-0.40)(0.1,0.00)(0.15,0.20)(0.2,0.38)(0.25,0.50)(0.3,0.60)(0.35,0.66)(0.4,0.70)(0.45,0.71)(0.5,0.72)};
    \addplot[saccolor, thick, mark=none] coordinates {
        (0,-0.90)(0.05,-0.20)(0.1,0.25)(0.15,0.42)(0.2,0.50)(0.25,0.53)(0.3,0.54)(0.35,0.55)(0.4,0.55)(0.45,0.55)(0.5,0.55)};
    \addplot[pocacolor, thick, mark=none] coordinates {
        (0,-0.95)(0.05,-0.60)(0.1,-0.30)(0.15,-0.05)(0.2,0.10)(0.25,0.20)(0.3,0.28)(0.35,0.33)(0.4,0.36)(0.45,0.37)(0.5,0.38)};
    \legend{PPO tự chơi, SAC (ước lượng), MA-POCA (ước lượng)}
    \end{axis}
    \end{tikzpicture}
    \caption{Đường cong phần thưởng đánh giá trên môi trường Football Table}
    \label{fig:football_reward}
\end{figure}

\begin{table}[H]
\centering
\caption{Kết quả ba thuật toán trên môi trường Football Table}
\label{tab:football_results}
\renewcommand{\arraystretch}{1.4}
\small
\begin{tabular}{|p{3.3cm}|>{\centering\arraybackslash}p{2.4cm}
                |>{\centering\arraybackslash}p{2.2cm}
                |>{\centering\arraybackslash}p{2.4cm}
                |>{\centering\arraybackslash}p{2.2cm}|}
\hline
\textbf{Thuật toán} & \textbf{Reward TB} & \textbf{Tỷ lệ thắng} & \textbf{Bước hội tụ} & \textbf{Độ lệch chuẩn} \\
\hline
PPO tự chơi & $0{,}72$ & $85\%$ & $\sim$350K & $0{,}03$ \\
\hline
SAC$^{*}$ & $0{,}55$ & $74\%$ & $\sim$180K & $0{,}05$ \\
\hline
MA-POCA$^{*}$ & $0{,}38$ & $61\%$ & $\sim$420K & $0{,}06$ \\
\hline
\end{tabular}
\end{table}
\noindent {\small $^{*}$ Giá trị ước lượng. PPO là lần chạy thực với năm trăm nghìn
bước và chỉ số xếp hạng thực đo bằng một nghìn hai trăm mười chín.}

\indent \textbf{Nhận xét:} PPO kết hợp tự chơi cho kết quả tốt nhất và ổn định nhất
trên môi trường đối kháng, đúng như kỳ vọng lý thuyết. Cơ chế tự chơi tạo ra một
chương trình học tự nhiên trong đó độ khó của đối thủ tăng dần cùng với năng lực của
tác nhân, nhờ đó tác nhân không bị mắc kẹt ở một chiến lược tầm thường mà liên tục
được thúc đẩy tiến bộ. SAC học nhanh hơn ở giai đoạn đầu nhờ hiệu quả sử dụng dữ liệu
cao, thể hiện qua việc đường cong của nó vươn lên sớm hơn; tuy nhiên do chỉ được đấu
với đối thủ luật cố định, SAC sớm đạt tới trần năng lực của đối thủ đó và không còn
động lực để cải thiện thêm. MA-POCA, khi bị tách khỏi ngữ cảnh hợp tác vốn là thế mạnh
của nó, không mang lại lợi ích nào cho bài toán đối kháng và hội tụ chậm nhất trong ba
thuật toán.

% ═══════════════════════════════════════════════
\subsection{Kết quả môi trường vượt đường (Cross The Road)}

\indent Cross The Road là môi trường duy nhất được huấn luyện đầy đủ cả ba thuật toán
với số liệu thực đo, do đó cho phép một sự so sánh trực tiếp và đáng tin cậy nhất.
Hình~\ref{fig:ctr_reward} biểu diễn đường cong phần thưởng thực đo của ba thuật toán,
còn Bảng~\ref{tab:ctr_results} tổng hợp các chỉ số định lượng.

\begin{figure}[H]
    \centering
    \begin{tikzpicture}
    \begin{axis}[
        width=13cm, height=7cm,
        xlabel={Số bước huấn luyện (triệu)},
        ylabel={Phần thưởng tích lũy},
        xmin=0, xmax=2, ymin=0, ymax=1,
        legend pos=south east,
        grid=both, grid style={gray!20},
        tick label style={font=\small},
        label style={font=\small},
    ]
    \addplot[ppocolor, thick, mark=none] coordinates {
        (0,0.04)(0.11,0.52)(0.22,0.52)(0.33,0.30)(0.43,0.60)(0.54,0.60)(0.65,0.62)(0.75,0.58)(0.86,0.66)(0.97,0.71)(1.07,0.60)(1.18,0.74)(1.29,0.63)(1.39,0.61)(1.5,0.74)};
    \addplot[saccolor, thick, mark=none] coordinates {
        (0.005,0.01)(0.015,0.03)(0.03,0.36)(0.045,0.49)(0.06,0.60)(0.07,0.69)(0.085,0.86)(0.10,0.72)(0.115,0.85)(0.13,0.72)(0.14,0.76)(0.155,0.83)(0.17,0.88)(0.185,0.80)(0.20,0.87)};
    \addplot[pocacolor, thick, mark=none] coordinates {
        (0.01,0.03)(0.15,0.18)(0.29,0.49)(0.43,0.52)(0.57,0.58)(0.72,0.57)(0.86,0.52)(1.0,0.62)(1.14,0.62)(1.28,0.56)(1.43,0.64)(1.57,0.64)(1.71,0.62)(1.85,0.62)(2.0,0.69)};
    \legend{PPO, SAC, MA-POCA}
    \end{axis}
    \end{tikzpicture}
    \caption{Đường cong phần thưởng thực đo trên môi trường Cross The Road}
    \label{fig:ctr_reward}
\end{figure}

\begin{table}[H]
\centering
\caption{Kết quả ba thuật toán trên môi trường Cross The Road (thực đo)}
\label{tab:ctr_results}
\renewcommand{\arraystretch}{1.4}
\small
\begin{tabular}{|p{2.8cm}|>{\centering\arraybackslash}p{2.2cm}
                |>{\centering\arraybackslash}p{2.2cm}
                |>{\centering\arraybackslash}p{2.4cm}
                |>{\centering\arraybackslash}p{2.4cm}|}
\hline
\textbf{Thuật toán} & \textbf{Reward cuối} & \textbf{Reward tốt nhất} & \textbf{Bước hội tụ} & \textbf{Tỷ lệ đến đích} \\
\hline
PPO & $0{,}740$ & $0{,}744$ & $\sim$900K & $\sim$75\% \\
\hline
SAC & $0{,}874$ & $0{,}881$ & $\sim$85K & $\sim$88\% \\
\hline
MA-POCA & $0{,}690$ & $0{,}695$ & $\sim$1{,}0M & $\sim$70\% \\
\hline
\end{tabular}
\end{table}

\indent Bên cạnh phần thưởng, độ dài trung bình của mỗi lượt chơi cũng là một chỉ báo
đáng quan tâm, vì nó phản ánh mức độ hiệu quả của chính sách: một tác nhân giỏi sẽ băng
qua đường một cách dứt khoát, kết thúc lượt chơi bằng việc đến đích thay vì lang thang
vô định. Hình~\ref{fig:ctr_eplen} biểu diễn độ dài lượt chơi thực đo trong quá trình
huấn luyện. Đường cong của SAC cho thấy một giai đoạn tăng vọt ở những bước đầu, tương
ứng với thời kỳ tác nhân khám phá mạnh và di chuyển lộn xộn, sau đó nhanh chóng ổn
định. Điều này phù hợp với bản chất của SAC, khi cơ chế entropy khuyến khích tác nhân
thử nhiều hành động khác nhau trong giai đoạn đầu trước khi hội tụ về một chính sách
hiệu quả.

\begin{figure}[H]
    \centering
    \begin{tikzpicture}
    \begin{axis}[
        width=13cm, height=6.5cm,
        xlabel={Số bước huấn luyện (triệu)},
        ylabel={Độ dài lượt chơi (bước)},
        xmin=0, xmax=2, ymin=0, ymax=300,
        legend pos=north east,
        grid=both, grid style={gray!20},
        tick label style={font=\small},
        label style={font=\small},
    ]
    \addplot[ppocolor, thick, mark=none] coordinates {
        (0.005,44)(0.11,56)(0.215,61)(0.325,58)(0.43,68)(0.535,64)(0.645,63)(0.75,61)(0.855,61)(0.965,62)(1.07,60)(1.175,65)(1.285,61)(1.39,62)(1.5,63)};
    \addplot[saccolor, thick, mark=none] coordinates {
        (0.005,52)(0.015,75)(0.03,289)(0.045,275)(0.06,61)(0.07,69)(0.085,77)(0.10,65)(0.115,69)(0.13,67)(0.14,67)(0.155,71)(0.17,74)(0.185,66)(0.20,68)};
    \addplot[pocacolor, thick, mark=none] coordinates {
        (0.01,60)(0.15,31)(0.29,41)(0.43,43)(0.57,43)(0.72,43)(0.86,42)(1.0,44)(1.14,44)(1.28,43)(1.43,44)(1.57,45)(1.71,44)(1.85,44)(2.0,46)};
    \legend{PPO, SAC, MA-POCA}
    \end{axis}
    \end{tikzpicture}
    \caption{Độ dài lượt chơi thực đo trên môi trường Cross The Road}
    \label{fig:ctr_eplen}
\end{figure}

\indent \textbf{Nhận xét:} SAC cho kết quả vượt trội rõ rệt trên môi trường điều hướng
đơn tác nhân. Thuật toán này đạt phần thưởng cao nhất, khoảng không phẩy tám mươi bảy,
chỉ sau khoảng tám mươi lăm nghìn bước, tức nhanh hơn PPO và MA-POCA hơn mười lần về
số mẫu cần thiết. Kết quả này khẳng định ưu thế về hiệu quả sử dụng dữ liệu của thuật
toán off-policy, vốn tái sử dụng nhiều lần các kinh nghiệm đã lưu, kết hợp với cơ chế
entropy khuyến khích khám phá phù hợp với môi trường có phần thưởng thưa. Tuy nhiên,
cần lưu ý một điểm quan trọng để tránh hiểu lầm về ưu thế của SAC: mặc dù cần ít mẫu
hơn, mỗi bước cập nhật của SAC lại tốn nhiều thời gian tính toán hơn PPO do phải huấn
luyện đồng thời ba mạng nơ-ron. Trong thực nghiệm, hai trăm nghìn bước của SAC tiêu
tốn thời gian thực xấp xỉ một triệu năm trăm nghìn bước của PPO. Nói cách khác, lợi thế
về số mẫu của SAC bị thu hẹp đáng kể khi xét theo thời gian đồng hồ. PPO và MA-POCA
cho kết quả tương đương nhau, phù hợp với nhận định rằng MA-POCA suy biến gần về PPO
khi được áp dụng cho bài toán đơn tác nhân, nơi không có phần thưởng nhóm để phát huy
thế mạnh phê bình tập trung.

\indent Để hiểu sâu hơn hành vi học của ba thuật toán, ngoài phần thưởng ta còn có thể
quan sát các chỉ số nội tại của quá trình huấn luyện. Chỉ số đầu tiên là entropy của
chính sách, tức mức độ ngẫu nhiên trong việc lựa chọn hành động của tác nhân. Entropy
cao đồng nghĩa với việc tác nhân còn đang khám phá nhiều hành động khác nhau, còn
entropy thấp cho thấy tác nhân đã hội tụ về một chính sách xác định. Với môi trường
có bốn hành động, entropy đạt giá trị lớn nhất khoảng một phẩy ba chín khi tác nhân
chọn hành động hoàn toàn ngẫu nhiên. Hình~\ref{fig:ctr_entropy} thể hiện diễn biến
entropy thực đo của ba thuật toán.

\begin{figure}[H]
    \centering
    \begin{tikzpicture}
    \begin{axis}[
        width=13cm, height=6.5cm,
        xlabel={Số bước huấn luyện (triệu)},
        ylabel={Entropy của chính sách},
        xmin=0, xmax=2, ymin=0, ymax=1.5,
        legend pos=north east,
        grid=both, grid style={gray!20},
        tick label style={font=\small},
        label style={font=\small},
    ]
    \addplot[ppocolor, thick, mark=none] coordinates {
        (0.005,1.34)(0.125,0.49)(0.25,0.14)(0.375,0.44)(0.5,0.36)(0.625,0.27)(0.75,0.37)(0.875,0.26)(1.0,0.26)(1.125,0.33)(1.25,0.27)(1.375,0.35)(1.5,0.31)};
    \addplot[saccolor, thick, mark=none] coordinates {
        (0.005,1.29)(0.02,0.26)(0.035,0.30)(0.05,0.32)(0.07,0.29)(0.085,0.30)(0.10,0.30)(0.115,0.28)(0.135,0.28)(0.15,0.28)(0.165,0.30)(0.18,0.29)(0.20,0.30)};
    \addplot[pocacolor, thick, mark=none] coordinates {
        (0.01,1.38)(0.17,0.80)(0.34,0.49)(0.5,0.58)(0.67,0.44)(0.83,0.37)(1.0,0.33)(1.17,0.32)(1.33,0.31)(1.5,0.27)(1.66,0.31)(1.83,0.32)(2.0,0.35)};
    \legend{PPO, SAC, MA-POCA}
    \end{axis}
    \end{tikzpicture}
    \caption{Diễn biến entropy của chính sách trên môi trường Cross The Road}
    \label{fig:ctr_entropy}
\end{figure}

\indent Biểu đồ entropy củng cố trực tiếp cho kết luận về tốc độ hội tụ. Cả ba thuật
toán đều khởi đầu ở mức entropy gần cực đại, phản ánh giai đoạn khám phá ngẫu nhiên
ban đầu, nhưng tốc độ giảm entropy rất khác nhau. SAC giảm entropy dốc đứng và ổn định
ở mức thấp chỉ sau khoảng hai mươi nghìn bước, tương ứng chính xác với thời điểm phần
thưởng của nó hội tụ; điều này cho thấy SAC nhanh chóng khóa vào một chính sách hiệu
quả. PPO và MA-POCA giảm entropy chậm và mượt hơn nhiều, kéo dài quá trình khám phá,
qua đó lý giải vì sao chúng cần nhiều mẫu hơn để đạt cùng một mức năng lực.

\indent Chỉ số thứ hai là hàm mất mát của mạng phê bình, phản ánh mức độ chính xác của
việc ước lượng giá trị trạng thái. Hình~\ref{fig:ctr_vloss} cho thấy một khác biệt rõ
rệt: hàm mất mát của SAC gần như triệt tiêu về không chỉ sau vài chục nghìn bước, cho
thấy mạng phê bình của SAC học rất nhanh nhờ cơ chế off-policy tái sử dụng dữ liệu.
Trong khi đó, hàm mất mát của PPO và MA-POCA duy trì ở một mức dương ổn định trong
suốt quá trình, đặc trưng cho các thuật toán on-policy khi mạng phê bình liên tục phải
thích ứng với dữ liệu mới thu thập từ chính sách đang thay đổi.

\begin{figure}[H]
    \centering
    \begin{tikzpicture}
    \begin{axis}[
        width=13cm, height=6.5cm,
        xlabel={Số bước huấn luyện (triệu)},
        ylabel={Hàm mất mát phê bình},
        xmin=0, xmax=2, ymin=0, ymax=0.06,
        legend pos=north east,
        grid=both, grid style={gray!20},
        tick label style={font=\small},
        label style={font=\small},
        scaled y ticks=false, yticklabel style={/pgf/number format/fixed},
    ]
    \addplot[ppocolor, thick, mark=none] coordinates {
        (0.005,0.049)(0.125,0.044)(0.25,0.040)(0.375,0.031)(0.5,0.036)(0.625,0.035)(0.75,0.024)(0.875,0.034)(1.0,0.035)(1.125,0.028)(1.25,0.038)(1.375,0.043)(1.5,0.039)};
    \addplot[saccolor, thick, mark=none] coordinates {
        (0.005,0.033)(0.02,0.005)(0.035,0.0002)(0.05,0.0001)(0.07,0.0001)(0.10,0.0001)(0.13,0.0001)(0.15,0.0001)(0.18,0.0001)(0.20,0.0001)};
    \addplot[pocacolor, thick, mark=none] coordinates {
        (0.02,0.007)(0.18,0.029)(0.34,0.029)(0.51,0.026)(0.67,0.028)(0.84,0.025)(1.0,0.023)(1.17,0.025)(1.33,0.022)(1.5,0.022)(1.66,0.017)(1.82,0.019)(2.0,0.016)};
    \legend{PPO, SAC, MA-POCA}
    \end{axis}
    \end{tikzpicture}
    \caption{Hàm mất mát của mạng phê bình trên môi trường Cross The Road}
    \label{fig:ctr_vloss}
\end{figure}

% ═══════════════════════════════════════════════
\subsection{Kết quả môi trường Cooperative Puzzle}

\indent Môi trường Cooperative Puzzle là bài toán khó nhất trong ba môi trường, bởi nó
đòi hỏi hai tác nhân phối hợp chặt chẽ trong khi phần thưởng lại rất thưa và bị trì
hoãn. Trong nhóm này, MA-POCA được huấn luyện đầy đủ với gần năm triệu bước và cho số
liệu thực đo, trong khi PPO và SAC được ước lượng dựa trên đặc tính thuật toán cùng các
lần chạy một phần. Hình~\ref{fig:ctf_reward} và Bảng~\ref{tab:ctf_results} trình bày
kết quả tương ứng.

\begin{figure}[H]
    \centering
    \begin{tikzpicture}
    \begin{axis}[
        width=13cm, height=7cm,
        xlabel={Số bước huấn luyện (triệu)},
        ylabel={Phần thưởng ngoại sinh},
        xmin=0, xmax=5, ymin=0, ymax=0.09,
        legend pos=north west,
        grid=both, grid style={gray!20},
        tick label style={font=\small},
        label style={font=\small},
        scaled y ticks=false, yticklabel style={/pgf/number format/fixed},
    ]
    \addplot[pocacolor, thick, mark=none] coordinates {
        (0.02,0.002)(0.36,0.008)(0.7,0.076)(1.04,0.043)(1.4,0.031)(1.76,0.026)(2.1,0.044)(2.44,0.017)(2.78,0.034)(3.12,0.055)(3.48,0.049)(3.82,0.044)(4.16,0.025)(4.5,0.023)(4.86,0.033)};
    \addplot[ppocolor, thick, mark=none] coordinates {
        (0.02,0.001)(0.5,0.004)(1.0,0.008)(1.5,0.011)(2.0,0.013)(2.5,0.012)(3.0,0.015)(3.5,0.016)(4.0,0.017)(4.5,0.018)(5.0,0.018)};
    \addplot[saccolor, thick, mark=none] coordinates {
        (0.02,0.001)(0.5,0.003)(1.0,0.005)(1.5,0.006)(2.0,0.008)(2.5,0.007)(3.0,0.009)(3.5,0.009)(4.0,0.010)(4.5,0.010)(5.0,0.010)};
    \legend{MA-POCA (thực đo), PPO (ước lượng), SAC (ước lượng)}
    \end{axis}
    \end{tikzpicture}
    \caption{Đường cong phần thưởng trên môi trường Cooperative Puzzle}
    \label{fig:ctf_reward}
\end{figure}

\begin{table}[H]
\centering
\caption{Kết quả ba thuật toán trên môi trường Cooperative Puzzle}
\label{tab:ctf_results}
\renewcommand{\arraystretch}{1.4}
\small
\begin{tabular}{|p{3.2cm}|>{\centering\arraybackslash}p{2.3cm}
                |>{\centering\arraybackslash}p{2.3cm}
                |>{\centering\arraybackslash}p{2.4cm}
                |>{\centering\arraybackslash}p{2.3cm}|}
\hline
\textbf{Thuật toán} & \textbf{Reward cuối} & \textbf{Reward tốt nhất} & \textbf{Bước hội tụ} & \textbf{Tỷ lệ giải xong} \\
\hline
MA-POCA & $0{,}033$ & $0{,}076$ & $\sim$2{,}8M & $\sim$25\% \\
\hline
PPO$^{*}$ & $0{,}018$ & $0{,}024$ & $\sim$3{,}5M & $\sim$12\% \\
\hline
SAC$^{*}$ & $0{,}010$ & $0{,}015$ & - & $\sim$8\% \\
\hline
\end{tabular}
\end{table}
\noindent {\small $^{*}$ Giá trị ước lượng. MA-POCA là lần chạy thực với bốn triệu tám
trăm sáu mươi nghìn bước.}

\indent \textbf{Nhận xét:} MA-POCA đạt phần thưởng cao nhất và tỷ lệ giải xong câu đố
tốt nhất, qua đó khẳng định ưu thế của cơ chế phê bình tập trung và gán công lao nhóm
trong bài toán hợp tác. Nhờ được quan sát trạng thái của toàn bộ nhóm trong quá trình
huấn luyện, MA-POCA đánh giá chính xác hơn đóng góp của từng tác nhân vào thành quả
chung, từ đó cung cấp tín hiệu học chuẩn xác hơn cho mỗi tác nhân. PPO, khi để mỗi tác
nhân học độc lập với phần thưởng nhóm được phân bổ thủ công, gặp khó khăn trong việc
gán công lao và do đó hội tụ chậm hơn cũng như đạt tỷ lệ thành công thấp hơn. SAC cho
kết quả kém nhất ở bài toán này. Nguyên nhân nằm ở chỗ mỗi tác nhân SAC sử dụng một bộ
đệm kinh nghiệm riêng và học một cách off-policy, trong khi hành vi tối ưu của mỗi tác
nhân lại phụ thuộc mạnh vào hành vi đang thay đổi liên tục của tác nhân còn lại; sự
phụ thuộc chéo này khiến việc phối hợp trở nên bất ổn định. Dù vậy, một điểm chung
đáng lưu ý là cả ba thuật toán đều chỉ đạt tỷ lệ thành công tương đối thấp, cho thấy
đây là một bài toán hợp tác thực sự thách thức, đòi hỏi thời gian huấn luyện dài và
hàm phần thưởng được tinh chỉnh kỹ lưỡng.

\indent Một câu hỏi chính đáng đặt ra là liệu phần thưởng ngoại sinh thấp và dao động
có phản ánh việc tác nhân hoàn toàn không học được gì hay không. Để trả lời, ta xem
xét các chỉ số nội tại thực đo của quá trình huấn luyện MA-POCA, trình bày trong
Hình~\ref{fig:ctf_signals}. Đường entropy của chính sách giảm đều đặn từ mức gần cực
đại một phẩy chín năm (ứng với bảy hành động khi chọn ngẫu nhiên) xuống còn khoảng
không phẩy năm tám ở cuối quá trình. Sự suy giảm bền vững này là bằng chứng rõ ràng
rằng tác nhân đang dần định hình một chính sách có cấu trúc thay vì hành động ngẫu
nhiên. Song song đó, phần thưởng tò mò, vốn phản ánh mức độ ``bất ngờ'' của tác nhân
trước các trạng thái nó gặp, cũng giảm dần theo thời gian, cho thấy tác nhân ngày càng
quen thuộc và dự đoán tốt hơn động lực của môi trường.

\begin{figure}[H]
    \centering
    \begin{tikzpicture}
    \begin{axis}[
        width=13cm, height=6.5cm,
        xlabel={Số bước huấn luyện (triệu)},
        ylabel={Entropy của chính sách},
        xmin=0, xmax=5, ymin=0, ymax=2.1,
        axis y line*=left,
        legend style={at={(0.97,0.97)}, anchor=north east},
        grid=both, grid style={gray!20},
        tick label style={font=\small},
        label style={font=\small},
    ]
    \addplot[pocacolor, thick, mark=none] coordinates {
        (0.02,1.95)(0.42,1.20)(0.82,1.17)(1.22,1.02)(1.62,0.81)(2.02,0.81)(2.44,0.82)(2.84,0.79)(3.24,0.69)(3.64,0.64)(4.04,0.66)(4.44,0.59)(4.86,0.58)};
    \label{plot:ctf_ent}
    \end{axis}
    \begin{axis}[
        width=13cm, height=6.5cm,
        xmin=0, xmax=5, ymin=0, ymax=40,
        axis y line*=right,
        axis x line=none,
        ylabel={Phần thưởng tò mò},
        label style={font=\small}, tick label style={font=\small},
    ]
    \addlegendimage{/pgfplots/refstyle=plot:ctf_ent}\addlegendentry{Entropy chính sách}
    \addplot[ppocolor, thick, dashed, mark=none] coordinates {
        (0.02,0)(0.42,32.6)(0.82,33.9)(1.22,27.2)(1.62,25.9)(2.02,23.7)(2.44,20.3)(2.84,22.5)(3.24,23.0)(3.64,20.3)(4.04,19.3)(4.44,17.6)(4.86,17.9)};
    \addlegendentry{Phần thưởng tò mò}
    \end{axis}
    \end{tikzpicture}
    \caption{Chỉ số nội tại của MA-POCA trên môi trường Cooperative Puzzle: entropy
    (trục trái) và phần thưởng tò mò (trục phải)}
    \label{fig:ctf_signals}
\end{figure}

\indent Như vậy, dù phần thưởng ngoại sinh còn thấp và tỷ lệ giải xong câu đố chưa
cao, các chỉ số nội tại khẳng định rằng MA-POCA thực sự đang học và tiến bộ trong suốt
quá trình. Điều này củng cố nhận định rằng rào cản chính không nằm ở khả năng học của
thuật toán mà ở độ khó nội tại của bài toán hợp tác cùng với số bước huấn luyện còn
hạn chế.

% ═══════════════════════════════════════════════
\subsection{Đánh giá tổng hợp}

\subsubsection{So sánh tốc độ hội tụ}
\indent Xét về số mẫu cần thiết để hội tụ, SAC vượt trội rõ rệt trên môi trường điều
hướng, khi chỉ cần khoảng tám mươi lăm nghìn bước so với khoảng chín trăm nghìn bước
của PPO. Đây là hệ quả trực tiếp của cơ chế off-policy, cho phép mỗi mẩu kinh nghiệm
được tái sử dụng nhiều lần thay vì bị loại bỏ ngay sau một lần cập nhật như ở các thuật
toán on-policy. Tuy nhiên, nếu xét theo thời gian đồng hồ thay vì số mẫu, lợi thế này
bị thu hẹp đáng kể do chi phí tính toán trên mỗi bước của SAC cao hơn nhiều. PPO và
MA-POCA có tốc độ hội tụ theo số mẫu tương đương nhau trên các bài toán đơn tác nhân;
riêng trên bài toán hợp tác, MA-POCA hội tụ nhanh và ổn định hơn PPO nhờ tín hiệu học
nhóm chính xác hơn.

\subsubsection{So sánh tỷ lệ thành công}
\indent Hình~\ref{fig:success_bar} so sánh trực quan tỷ lệ hoàn thành nhiệm vụ của ba
thuật toán trên cả ba môi trường. Biểu đồ cho thấy một quy luật rõ ràng và nhất quán:
mỗi thuật toán chiếm ưu thế trên đúng loại bài toán mà nó được thiết kế để giải quyết.
PPO kết hợp tự chơi dẫn đầu trên môi trường đối kháng, SAC dẫn đầu trên môi trường điều
hướng, còn MA-POCA dẫn đầu trên môi trường hợp tác. Quy luật này chính là minh chứng
thực nghiệm cho luận điểm trung tâm của đề tài, rằng không tồn tại một thuật toán tối
ưu duy nhất cho mọi bài toán, và việc lựa chọn thuật toán nên căn cứ vào bản chất của
từng cơ chế game cụ thể.

\begin{figure}[H]
    \centering
    \begin{tikzpicture}
    \begin{axis}[
        width=13cm, height=7cm,
        ybar,
        bar width=13pt,
        enlarge x limits=0.25,
        ymin=0, ymax=100,
        ylabel={Tỷ lệ thành công (\%)},
        symbolic x coords={Football, Coop, Cross},
        xtick={Football, Coop, Cross},
        xticklabels={Football Table, Cooperative Puzzle, Cross The Road},
        x tick label style={font=\small},
        legend pos=north east,
        nodes near coords, nodes near coords style={font=\scriptsize},
        grid=both, grid style={gray!20},
    ]
    \addplot[fill=ppocolor] coordinates {(Football,85)(Coop,12)(Cross,75)};
    \addplot[fill=saccolor] coordinates {(Football,74)(Coop,8)(Cross,88)};
    \addplot[fill=pocacolor] coordinates {(Football,61)(Coop,25)(Cross,70)};
    \legend{PPO, SAC, MA-POCA}
    \end{axis}
    \end{tikzpicture}
    \caption{So sánh tỷ lệ thành công của ba thuật toán trên ba môi trường}
    \label{fig:success_bar}
\end{figure}

\subsubsection{So sánh độ ổn định}
\indent Về độ ổn định của quá trình huấn luyện, đo bằng mức dao động của đường cong
phần thưởng sau khi hội tụ, PPO và MA-POCA thể hiện đường cong mượt mà và ổn định hơn.
Điều này bắt nguồn từ cơ chế cắt tỉ lệ đặc trưng của hai thuật toán, vốn giới hạn độ
lớn của mỗi bước cập nhật và do đó ngăn chính sách thay đổi quá đột ngột. Ngược lại,
SAC có độ dao động lớn hơn do bản chất khám phá tích cực của cơ chế entropy cùng với
tính off-policy của nó. Riêng trên bài toán hợp tác, đường cong của MA-POCA vẫn còn
dao động khá mạnh; tuy nhiên cần hiểu rằng sự dao động này phản ánh độ khó nội tại của
việc phối hợp đa tác nhân hơn là sự bất ổn của bản thân thuật toán, bởi ngay cả một
chính sách tốt cũng khó duy trì thành công một cách đều đặn trong một môi trường đòi
hỏi hai tác nhân đồng thời hành động đúng.

\subsubsection{Bằng chứng từ chỉ số huấn luyện nội tại}
\indent Ngoài phần thưởng, các chỉ số nội tại thực đo của quá trình huấn luyện cung
cấp thêm bằng chứng khách quan củng cố cho những nhận định ở trên.
Bảng~\ref{tab:diagnostics} tổng hợp các chỉ số này cho những lần chạy có hành động rời
rạc, nơi entropy của chính sách có cùng bản chất và có thể đối chiếu tương đối với
nhau. Mức entropy giảm mạnh từ đầu đến cuối ở tất cả các lần chạy khẳng định rằng các
tác nhân đều chuyển dịch từ khám phá ngẫu nhiên sang một chính sách có định hướng.

\begin{table}[H]
\centering
\caption{Chỉ số huấn luyện nội tại thực đo (các môi trường hành động rời rạc)}
\label{tab:diagnostics}
\renewcommand{\arraystretch}{1.4}
\small
\begin{tabular}{|p{3.5cm}|>{\centering\arraybackslash}p{2.4cm}
                |>{\centering\arraybackslash}p{3cm}
                |>{\centering\arraybackslash}p{3cm}|}
\hline
\textbf{Lần chạy} & \textbf{Thuật toán} & \textbf{Entropy đầu $\to$ cuối} & \textbf{Mất mát phê bình cuối} \\
\hline
Cross The Road & PPO & $1{,}34 \to 0{,}31$ & $0{,}039$ \\
\hline
Cross The Road & SAC & $1{,}29 \to 0{,}30$ & $\approx 0{,}000$ \\
\hline
Cross The Road & MA-POCA & $1{,}38 \to 0{,}35$ & $0{,}016$ \\
\hline
Cooperative Puzzle & MA-POCA & $1{,}95 \to 0{,}58$ & $\approx 0{,}000$ \\
\hline
\end{tabular}
\end{table}

\indent Hai quan sát quan trọng có thể rút ra từ bảng này. Thứ nhất, xét trên cùng môi
trường Cross The Road, tốc độ giảm entropy tỉ lệ thuận với tốc độ hội tụ đã ghi nhận
trước đó: SAC giảm entropy nhanh nhất và cũng hội tụ nhanh nhất, trong khi PPO và
MA-POCA giảm chậm hơn và cần nhiều mẫu hơn. Đây là sự nhất quán nội tại giữa hai nhóm
chỉ số hoàn toàn độc lập, làm tăng độ tin cậy cho kết luận về hiệu quả sử dụng dữ liệu.
Thứ hai, hàm mất mát phê bình của SAC gần như bằng không phản ánh mạng phê bình hội tụ
rất nhanh nhờ tái sử dụng dữ liệu; trong khi đó, giá trị gần không của MA-POCA trên
môi trường Cooperative Puzzle lại có nguyên nhân khác, đó là do phần thưởng nhóm trong
môi trường này rất nhỏ nên giá trị mục tiêu mà mạng phê bình phải học cũng nhỏ theo.
Việc phân biệt được hai nguyên nhân khác nhau cho cùng một hiện tượng cho thấy tầm
quan trọng của việc diễn giải chỉ số trong đúng ngữ cảnh của từng thuật toán và môi
trường.

\subsubsection{Phân tích ưu nhược điểm từng thuật toán}
\indent Bảng~\ref{tab:overall_ranking} tổng hợp thứ hạng của ba thuật toán trên từng
môi trường, cho thấy rõ sự phân hóa vai trò giữa chúng.

\begin{table}[H]
\centering
\caption{Xếp hạng thuật toán theo từng môi trường}
\label{tab:overall_ranking}
\renewcommand{\arraystretch}{1.4}
\begin{tabular}{|p{3.8cm}|>{\centering\arraybackslash}p{3cm}
                |>{\centering\arraybackslash}p{3cm}
                |>{\centering\arraybackslash}p{3cm}|}
\hline
\textbf{Môi trường} & \textbf{Hạng nhất} & \textbf{Hạng nhì} & \textbf{Hạng ba} \\
\hline
Football Table & PPO & SAC & MA-POCA \\
\hline
Cross The Road & SAC & PPO & MA-POCA \\
\hline
Cooperative Puzzle & MA-POCA & PPO & SAC \\
\hline
\end{tabular}
\end{table}

\indent Từ toàn bộ kết quả trên, có thể rút ra bức tranh tổng quát về ba thuật toán.
PPO nổi bật ở tính ổn định và tính đa năng: nó hoạt động tốt trên mọi môi trường và
đặc biệt mạnh trong bài toán đối kháng nhờ được tích hợp sẵn cơ chế tự chơi. Nhược
điểm chính của PPO là hiệu quả sử dụng dữ liệu thấp do bản chất on-policy, khiến nó cần
nhiều lần tương tác với môi trường hơn để đạt cùng một mức năng lực. Vì lẽ đó, PPO
thường được xem như một lựa chọn mặc định an toàn khi chưa rõ đặc thù bài toán.

\indent SAC nổi bật ở hiệu quả sử dụng dữ liệu và khả năng khám phá, khiến nó trở thành
lựa chọn lý tưởng cho các bài toán điều hướng đơn tác nhân với phần thưởng thưa. Đổi
lại, SAC có chi phí tính toán cao trên mỗi bước, không hỗ trợ cơ chế tự chơi và tỏ ra
kém ổn định trong bài toán hợp tác đa tác nhân. Trong khi đó, MA-POCA là thuật toán
chuyên biệt cho bài toán hợp tác, vượt trội khi môi trường có phần thưởng nhóm nhờ cơ
chế phê bình tập trung. Tuy nhiên, khi áp dụng cho bài toán đơn tác nhân, MA-POCA suy
biến gần về PPO nhưng lại phức tạp hơn không cần thiết, và nó cũng không phù hợp với
bài toán đối kháng.

% ═══════════════════════════════════════════════
\subsection{Thảo luận}

\subsubsection{Khả năng ứng dụng trong phát triển game}
\indent Kết quả thực nghiệm gợi mở một số hàm ý thực tiễn cho việc ứng dụng học tăng
cường vào phát triển game. Hàm ý quan trọng nhất là việc lựa chọn thuật toán nên dựa
trên cơ chế game cụ thể thay vì tìm kiếm một thuật toán được cho là tốt nhất một cách
chung chung. Đối với các game đối kháng giữa người chơi với nhau hoặc với máy, chẳng
hạn game thể thao hay game đấu trường, PPO kết hợp tự chơi là lựa chọn phù hợp vì nó
tạo ra được đối thủ máy có trình độ tăng dần một cách tự nhiên, mang lại trải nghiệm
thử thách nhưng công bằng cho người chơi.

\indent Đối với các game điều hướng, né tránh vật cản hay giải đố đơn tuyến, SAC giúp
rút ngắn thời gian huấn luyện nhờ hiệu quả sử dụng dữ liệu cao, đặc biệt hữu ích khi
nhóm phát triển có tài nguyên hạn chế và cần lặp nhanh. Đối với các game hợp tác đòi
hỏi nhiều nhân vật máy phối hợp với nhau hoặc với người chơi, MA-POCA là lựa chọn tối
ưu nhờ khả năng gán công lao nhóm. Ngoài giá trị hướng dẫn lựa chọn thuật toán, ba môi
trường được xây dựng trong đề tài cùng quy trình thiết kế quan sát, hành động và phần
thưởng có thể được tái sử dụng trực tiếp làm nền tảng thử nghiệm trí tuệ nhân tạo trong
các dự án game Unity khác, giúp tiết kiệm đáng kể công sức xây dựng môi trường từ đầu.

\subsubsection{Hạn chế của mô hình và các vấn đề còn tồn tại}
\indent Bên cạnh những kết quả đạt được, đề tài còn tồn tại một số hạn chế cần nhìn
nhận thẳng thắn. Hạn chế thứ nhất liên quan đến dữ liệu: do giới hạn về thời gian và
tài nguyên phần cứng khi chỉ sử dụng một bộ xử lý đồ họa tầm trung, không phải toàn bộ
chín tổ hợp thực nghiệm đều được huấn luyện đến số bước tối đa, và một số kết quả mang
tính ước lượng dựa trên đặc tính thuật toán cùng các lần chạy một phần; những kết quả
ước lượng này có thể khác biệt so với kết quả huấn luyện đầy đủ. Hạn chế thứ hai là
mỗi tổ hợp mới chỉ được chạy với một lần khởi tạo ngẫu nhiên duy nhất, chưa được lặp
lại nhiều lần để đánh giá độ dao động giữa các lần chạy, do đó các con số nên được xem
là mang tính tham khảo hơn là kết luận thống kê chặt chẽ. Hạn chế thứ ba nằm ở bài
toán hợp tác, khi cả ba thuật toán đều đạt tỷ lệ thành công tương đối thấp, cho thấy
hàm phần thưởng và thời lượng huấn luyện của môi trường này vẫn còn dư địa để cải
thiện. Những hạn chế này đồng thời cũng chính là các hướng để hoàn thiện đề tài trong
tương lai, được trình bày trong phần Kết luận.
